\section{Overview}
\label{sec:overview}

pyALDIC computes full-field displacement and strain from a sequence
of grayscale images of a deforming specimen. The typical workflow,
from left to right through the GUI, is:

\begin{enumerate}
  \item \textbf{Load images} (a folder of frames, first frame = reference).
  \item \textbf{Choose Workflow Type} --- tracking mode, solver, and
    reference-update policy. These choices are made \emph{before}
    drawing a Region of Interest because they determine which frames
    need one.
  \item \textbf{Draw a Region of Interest} on the required frame(s).
  \item \textbf{Set DIC parameters} --- subset size, node spacing,
    search range, mesh refinement.
  \item \textbf{Run the analysis.} Progress, Cancel, and console log
    are on the right panel.
  \item \textbf{View displacement results} live in the canvas.
  \item \textbf{Compute strain} in the Strain Post-Processing window
    (opens automatically when the run finishes).
  \item \textbf{Export} numerical data, images, or animations.
\end{enumerate}

\begin{notebox}
pyALDIC does \textbf{not} require you to pre-compute calibration,
stereo matching, or any 3D reconstruction. It is strictly 2D DIC
(one camera, planar specimen). If the specimen moves out of plane or
the illumination changes significantly between frames, results will
degrade and the GUI will warn you.
\end{notebox}

\subsection{Three-column layout}

The main window has three columns:

\begin{description}[leftmargin=1.2cm,style=nextline]
  \item[Left sidebar]
    Images panel, \emph{Workflow Type} section, \emph{Region of
    Interest} section, \emph{Parameters} section, and a collapsible
    \emph{Advanced} section.
  \item[Centre canvas]
    Zoom / pan viewport showing the current frame with overlays
    (ROI, mesh, displacement field).
  \item[Right sidebar]
    Run button, progress bar, field selector (disp\_u, disp\_v,
    etc.), visualization controls (colormap, range, opacity),
    physical units, and console log.
\end{description}

\subsection{Typical session times}

On a modern laptop (8-core CPU, 16 GB RAM):

\begin{center}
\begin{tabular}{lll}
  \toprule
  Image size & Grid spacing & Time per frame \\
  \midrule
  $256 \times 256$  & step 8   & $\sim$ 0.2 s \\
  $512 \times 512$  & step 8   & $\sim$ 1 s   \\
  $1024 \times 1024$ & step 4  & $\sim$ 3 s   \\
  \bottomrule
\end{tabular}
\end{center}

Times scale linearly with the number of mesh nodes. The
\emph{Local DIC} solver is roughly $3\times$ faster than
\emph{AL-DIC}.
